\begin{beginnerstatement}[Kurz erklärt: Vite und TypeScript]

Das Frontend ist die sichtbare und bedienbare Seite der Anwendung. JavaScript
ist eine Sprache für solche Anwendungen; TypeScript erweitert JavaScript um
Typangaben, mit denen bestimmte Fehler schon während der Entwicklung auffallen.
Vite ist das Werkzeug, das einen Entwicklungsserver startet und später den
Build erzeugt. Der Entwicklungsserver stellt die Anwendung beim Programmieren
lokal bereit, während der Build die auslieferbaren Dateien in
\path{frontend/dist} erstellt. Vitest führt automatisierte Tests für
Frontend-Code aus. Der HTTP-Client schickt Anfragen an ein Backend und verwendet
dazu eine URL, also dessen Adresse im Netzwerk. Im Template wird dieser Zugriff
nur für Profile mit Backend aktiviert.

\end{beginnerstatement}
